\documentclass[final,a4paper,12pt]{article}
\usepackage{notes}

\newcommand{\half}{\tfrac{1}{2}}

\title{Problem Set: Numerical Methods}
\author{Geophysics 3A}
\date{}

\begin{document}
\maketitle

\section{Background}

Cooling through oceanic seafloor is relatively well-characterized by a plate
cooling model. Heat flow patterns at the surface are more complicated as
hydrothermal circulation redistributes heat in the upper oceanic crust.
Sediments covering the oceanic crust cut off ventilated flow directly between
the igneous crust and the oceans. However, the igneous basement has sufficiently
high permeability that fluids can continue to circulate beneath the sediments.

In this practical, we will use a finite difference solution to the heat equation
to model heat flow through sediments and to the surface.

Modeling heat flow through sediments is complicated---not because the physics is
complex---but because the geometry is irregular. The depth to the
sediment--basement interface varies laterally, and sediment compaction alters
the physical properties of the sediments with depth.

\textbf{Important assumption:} Thermal conductivity and heat production vary
\textbf{only with depth} due to compaction, while sediment thickness varies
laterally. This combination produces a fully two-dimensional temperature and
heat flow field.


Because of these complexities, a numerical approach is required rather than an
analytic solution.

\bigskip
\section{Finite Difference Formulation}



We will construct a 2-D steady-state solution to the heat equation. In the
numerical grid:

\begin{itemize}
    \item Temperatures and depth are defined at \textbf{nodes}
    \item Thermal conductivity and heat production are defined at
    \textbf{cell centers}
\end{itemize}

Follow the steps below carefully.

\bigskip

\begin{enumerate}

\item \textbf{Governing equation}

\textbf{Task:} Write down the partial differential equation governing steady-state
heat conduction in two dimensions with spatially variable thermal conductivity
$k$ and heat production $A$.

\textbf{Hint:} Combine conservation of energy with Fourier's law.

\bigskip

\item \textbf{Boundary conditions}

The temperatures at the ocean--sediment interface and the sediment--basement
interface are fixed. However, the side boundaries require special treatment.

\begin{enumerate}

\item \textbf{Top and bottom boundaries}

\textbf{Task:} Write expressions for:
\begin{itemize}
    \item the surface temperature
    \item the basal temperature
\end{itemize}

\textbf{Note:} The basal temperature may vary laterally.

\medskip

\item \textbf{Side boundaries (insulating)}

We assume no heat can flow across the vertical sides of the model.

\textbf{Task:}
\begin{itemize}
    \item Write the mathematical condition for zero heat flux
    \item Express this condition in finite difference form
\end{itemize}

\end{enumerate}

\bigskip

\item \textbf{Finite difference equation}

\begin{figure}[h!]
\begin{center}
\begin{tikzpicture}[scale=2.5]

    % --- Draw grid ---
    \draw[gray!40] (-2,-2) grid (2,2);

    % --- Shade central cells ---
    \foreach \x/\z in {-0.5/-0.5, 0.5/-0.5, -0.5/0.5, 0.5/0.5} {
        \fill[blue!10] (\x-0.5,\z-0.5) rectangle (\x+0.5,\z+0.5);
    }

    % --- Cell-centred properties (with half indices) ---
    \node at (-0.5,-0.5) {\small 
        $\begin{aligned}
        k_{i-\half,j+\half} \\
        A_{i-\half,j+\half}
        \end{aligned}$
    };
    \node at (0.5,-0.5)  {\small 
        $\begin{aligned}
            k_{i+\half,j+\half} \\
            A_{i+\half,j+\half}
        \end{aligned}$
    };
    \node at (-0.5,0.5)  {\small 
        $\begin{aligned}
            k_{i-\half,j-\half} \\
            A_{i-\half,j-\half}
        \end{aligned}$
    };
    \node at (0.5,0.5)   {\small 
        $\begin{aligned}
            k_{i+\half,j-\half} \\
            A_{i+\half,j-\half}
        \end{aligned}$
    };

    % --- Nodes (temperatures) ---
    \fill (0,0) circle (1pt) node[above right] {$T_{i,j}$};
    \fill (1,0) circle (1pt) node[right] {$T_{i+1,j}$};
    \fill (-1,0) circle (1pt) node[left] {$T_{i-1,j}$};
    \fill (0,1) circle (1pt) node[above] {$T_{i,j-1}$};
    \fill (0,-1) circle (1pt) node[below] {$T_{i,j+1}$};

    % --- Stencil links ---
    \draw[dashed] (0,0) -- (1,0);
    \draw[dashed] (0,0) -- (-1,0);
    \draw[dashed] (0,0) -- (0,1);
    \draw[dashed] (0,0) -- (0,-1);

    % --- Grid spacing ---
    \draw[<->] (-2.2,0) -- (-2.2,1)
        node[midway,left] {$\Delta z$};

    \draw[<->] (-1,-2.2) -- (0,-2.2)
        node[midway,below] {$\Delta x$};

\end{tikzpicture}
\end{center}
\caption{Stencil for the 2-D thermal problem.  The temperatures are \emph{node centered} (defined at the nodes, circles).  The thermal properties are \emph{cell centered} (defined at the center of the cells, rectangles).  The cell positions are defined in terms of $(i,j)$ as the (column, row) positions associated with $(x,z)$.}
\label{fig:stencil}
\end{figure}


\textbf{Task:} Derive the finite difference equation for the temperature at node
$(i,j)$ using a central difference approximation.

\textbf{Important:}
\begin{itemize}
    \item Write the equation in terms of heat fluxes entering and leaving the node
    \item Use conductivities defined at the \textbf{interfaces} (up, down, left, right)
\end{itemize}

\textbf{Your answer should clearly show:}
\begin{itemize}
    \item contributions from all four neighbouring nodes
    \item the role of heat production at the node
\end{itemize}

\bigskip

\item \textbf{Heat production at nodes}

Heat production is defined at cell centres, but the finite difference equation
requires values at nodes.

\textbf{Task:}
\begin{enumerate}
    \item Derive an expression for heat production at node $(i,j)$ using the
    arithmetic mean of the four surrounding cells.
    \item \textbf{Then simplify:} Assume heat production varies only with depth
    (i.e.\ $A = A(z)$). What does your expression reduce to?
\end{enumerate}

\bigskip

\item \textbf{Thermal conductivity in four directions}

Thermal conductivity varies with depth due to compaction.

\textbf{Task:} Write expressions for conductivity in the four directions
(up, down, left, right) from node $(i,j)$.

\textbf{Hint:} Use the arithmetic mean of adjacent cell values in the direction
of heat flow.

\bigskip

\item \textbf{Porosity-controlled properties}

Sediment compaction reduces porosity with depth according to:
\begin{equation}
\phi = \phi_0 \exp(-z/\lambda)
\end{equation}

\bigskip

\textbf{Tasks:}

\begin{enumerate}
    \item Thermal conductivity depends on both the sediment matrix and pore fluid.
    Using a harmonic mean, show that:
    \[
    k = \left( \frac{1-\phi}{k_m} + \frac{\phi}{k_w} \right)^{-1}
    \]

    \item Heat production is associated with the solid matrix only. Show that:
    \[
    A = A_0 (1-\phi)
    \]

    \item \textbf{Explain qualitatively:}
    \begin{itemize}
        \item How does conductivity vary with depth?
        \item How does heat production vary with depth?
    \end{itemize}

\end{enumerate}

\bigskip

\item \textbf{Surface heat flow}

\textbf{Task:} Derive an expression for vertical surface heat flow using your
finite difference solution.

\textbf{Hint:} Use the temperature gradient between the uppermost nodes.

\bigskip

\item \textbf{Conceptual question (important)}

\textbf{Question:}

\emph{If thermal conductivity $k$ and heat production $A$ vary only with depth
(and not laterally), what causes the temperature and heat flow field to be
two-dimensional?}

\textbf{Answer clearly in words.}

\medskip

\textbf{Hint:} Think about the role of geometry.

\bigskip

\item \textbf{Sketching exercise}

Draw a cross section across a buried bathymetric high.

\textbf{Tasks:}
\begin{itemize}
    \item Draw two cases:
    \begin{enumerate}
        \item 1-D (uniform sediment thickness)
        \item 2-D (varying sediment thickness)
    \end{enumerate}
    \item On each sketch, show:
    \begin{itemize}
        \item isotherms
        \item heat flow directions
    \end{itemize}
\end{itemize}

\end{enumerate}

\bigskip
\section{Computational Exercise}

After completing the formulation, run \verb+iberia2d.py+.

The data are from heat flow measurements collected on the Iberian Abyssal Plain.
Sediment thickness is derived from seismic data, and thermal properties are
estimated from measurements and core samples.

The goal is to develop a 2-D thermal model that matches the observations.

\bigskip

\begin{enumerate}
\setcounter{enumi}{9}

\item \textbf{Model comparison}

The blue model is 1-D and the red model is 2-D.

\textbf{Tasks:}
\begin{itemize}
    \item Describe how well each model fits the observed data
    \item Identify where each model performs well or poorly
\end{itemize}

\bigskip

\item \textbf{3-D effects}

\textbf{Task:} Describe how heat flow over a 2-D ridge differs from that over a
buried seamount (a 3-D feature).

\textbf{Answer in words:}
\begin{itemize}
    \item How do heat flow patterns differ?
    \item Would including 3-D effects significantly reduce the misfit?
\end{itemize}

\bigskip

\item \textbf{Improving the model}

\textbf{Task:} If 3-D geometry does not fully explain the data, suggest ways to
improve the model.

\end{enumerate}

\end{document}

% \documentclass[final,a4paper,12pt]{article}
% \usepackage{notes}

% \title{Problem Set: Numerical Methods}
% \author{Geophysics 3A}
% \date{}

% \begin{document}
% \maketitle

% \section{Background}
% Cooling through oceanic seafloor is relatively well-characterized by a plate
% cooling model.  Heat flow patterns at the surface are more complicated as
% hydrothermal circulation redistributes the heat in the upper oceanic crust.
% Sediments covering the oceanic crust cut off ventilated flow directly between
% the igneous crust and oceans.  The igneous basement has a high enough
% permeability that fluids can continue to circulate
% vigorously.

% In this practical, we'll use a finite difference solution to the heat
% equation to model to model heat flow through the sediments and out of the
% surface.

% \section{Finite Differences}
% Modeling heat flow through the sediments is complicated---not because it is
% complex---by a variable depth to the sediment--basement interface.
% An additional complexity arises as a result of sediment compaction which alters
% the physical properties (thermal conductivity and heat production) of the
% sediment.  While neither complexity is difficult to model, they can make it
% considerably more difficult to develop an analytic solution.  Hence, a numerical
% approach is the best option for a solution.

% Where sediment cover is thick enough to blanket all basement highs, fluid flow
% across the crust--ocean interface is cutoff.  The sediments insulate the
% lithosphere, but fluid circulation within the igneous basement continues
% maintaining a constant temperature at the sediment--basement interface.  Bottom
% water temperatures at the ocean--sediment interface are also constant.  Thermal
% equilibrium is established within the sediments above.

% Let's develop 2-D steady-state temperature solution to solve for heat flow
% through the sediments.  Heat production and thermal conductivity will be
% assigned to the cells and temperature and depth to the nodes.  Follow the
% following steps below to derive a method for numerically computing heat flow:
% \newcounter{savenum}
% \begin{enumerate}
%     \item What PDE governs the system?  Remember, we are doing this for a 2-D
%         volume, incorporating a changing thermal conductivity and heat
%         production with depth.

%     \item Temperatures at the ocean--crust and sediment--basement interfaces are
%     simple since they are fixed, but the insulating conditions at the side
%     must be handled separately.  Write an expression for the temperatures on
%     either vertical boundary of the model.
%     \begin{enumerate}
%         \item What are the surface and basal boundary conditions? (Hint: they are
%         described above).

%         \item We will assume insulating conditions (zero-flux) on the vertical sides
%         of our boundary.  
%     \end{enumerate}

%     \item What is the central difference solution for $T_{i,j}$ at a location
%         $(z_i, x_j)$?  At the moment, use $k_{i,j}$ as the effective
%         conductivity and $A_{i,j}$ as the effective heat production at $(z_i,
%         x_j)$.

%     \item What is the average heat production at node $(i,j)$, estimated from
%         the weighted arithmetic mean of the surrounding cells?

%     \item For the thermal conductivity, we need to take into account the
%         anisotropy since the in conductivity varies vertically due to
%         compaction.  This means we need to know the thermal conductivity, for
%         heat flowing in four directions, up, down, left and right from each
%         node $(i,j)$.  Write an express for the thermal conductivity in each of the four
%         directions, using the arithmetic mean of the two conductivities in each
%         of the cells in the direction heat is flowing.

%     \item We'll assume porosity increases with depth in the sediments according
%         to the relationship
%         \begin{equation}
%             \phi = \phi_0 \exp\left( -z / \lambda \right)
%         \end{equation}
%         where $\lambda$ is the compaction length and $\phi_0$ is the surface
%         porosity.  Now write a relationship for the conductivity and heat
%         production with depth knowing the conductivity of the matrix is $k_m$
%         and the heat production of the matrix is $A_m$.  (You've did something
%         similar a couple assignments ago.)

%     \item Once we have a temperature model, we need to estimate the vertical
%         surface heat flow using our finite difference solution.
%     \setcounter{savenum}{\value{enumi}}

%     \item Before we move onto the computational portion of the practical, draw a
%         cross section across a 2-D bathymetric high.  Repeat the drawing so that
%         you have two copies to the left and right of each other.  Now draw the
%         expected pattern of heat flow and isotherms expected for the 1-D and 2-D
%         cases.
% \end{enumerate}

% Following the completion of your formulation, use the finite difference
% formulation in \verb+ssthermal.m+ to solve for the 2D heat flow.
% Background on the problem you will be solving: the data we are modeling comes
% from a set of heat flow data collected on the Iberian plain on the Atlantic
% seafloor.  Seismic data was collected which gives us an estimate of the sediment
% thickness.  Thermal conductivity is estimated either during measurement, or on
% sediment cores on the ship.  The goal of this modeling exercise is to develop a
% 2D thermal model which fit the data.  Initial modeling with a constant
% sediment-basement temperature did not fit the data properly so a simple linear
% increase in this basal temperature towards the continental shelf was
% incorporated.

% Make sure you have downloaded all the necessary codes and data files.

% Run \verb+iberia2d+.

% \begin{enumerate}\addtocounter{enumi}{\value{savenum}}
%     \item The blue model is 1-D and the red model is 2-D.  How well does each
%     model fit the observed data?

%     \item This model is 2-D, but bathymetry variations occur in 3-D.  Describe
%     (in words, no numbers/equations necessary) how the heat flow across a
%     2-D ridge would differ from a buried seamount.  Would accounting for 3-D
%     bathymetry variations account for the much of the misfit in the data?

%     \item If not all the data can be described by a move to 3-D, how might we
%     improve this model?
% \end{enumerate}

% \end{document}
